% GENERATED FILE. Edit canonical lessons, then run npm run generate:books.
% canonical-source-hash: fnv1a64:cdb43d2d0690fff3
% canonical-lessons: BN-C449-sneh, BN-C449-maya, BN-C449-ador, BN-C449-prem, BN-C449-bhokti-kora

\chapter{Affection}
\label{ch:bn-affection}
\hlchaptermodality{449}

\begin{chapteropening}
\textbf{By the end of this chapter:} \emph{I can talk about affection and love in a family.}

Complete the last lesson of chapter 449: I can talk about affection and love in a family.
\end{chapteropening}

\section[\texorpdfstring{\bn{স্নেহ} (sneh)}{স্নেহ (sneh)}]{\bn{স্নেহ} (sneh) — affection; tender love for someone younger}

\label{lesson:BN-C449-sneh}



\begin{quote}
Before the new one: say the Bengali for anxious with longing, then the Bengali for to cry out in despair.
\end{quote}

\subsection*{You'll want to know: \bn{স্নেহ}}
\textbf{\bn{স্নেহ}} — \emph{sneh} — \textquotedblleft{}affection; tender love for someone younger\textquotedblright{}.

\begin{quote}
\bn{দাদু} \bn{নাতিকে} \bn{স্নেহ} \bn{করেন।}
\end{quote}

\begin{quote}
\emph{dādu nātike sneh kôren.}
\end{quote}

\begin{quote}
Grandfather is very fond of his grandson.
\end{quote}

\bn{স্নেহ} runs downward, from an elder to someone younger; what runs back up is \bn{সম্মান} (taught), respect.

\begin{cousinweb}
\textbf{\bn{স্নেহ}} (\emph{sneh}) is a tatsama, Sanskrit \emph{sneha}, which Monier-Williams gives first as \textquotedblleft{}oiliness, smoothness\textquotedblright{} and then as \textquotedblleft{}tenderness, fondness\textquotedblright{}; Bengali keeps only the tenderness.
\end{cousinweb}

\begin{grammarlens}[title={What you've built}]
One more word for describing what you lived and felt.
\end{grammarlens}

\subsection*{Your turn}
\begin{itemize}
\raggedright
  \item \emph{Say it:} \emph{sneh}
  \item \emph{Say it:} \emph{sneh}, once more
  \item \emph{Say it:} \emph{dādu nātike sneh kôren.} — Grandfather is very fond of his grandson.
  \item \emph{Recall:} say the Bengali for anxious with longing, then the Bengali for to cry out in despair, then say \emph{sneh} again
\end{itemize}

\begin{tcolorbox}[breakable,colback=teal!4,colframe=teal!35!black,title={Before you move on}]
What does \bn{স্নেহ} mean? (\textquotedblleft{}Affection; tender love for someone younger\textquotedblright{}.) Say it once more.
\end{tcolorbox}
\section[\texorpdfstring{\bn{মায়া} (māyā)}{মায়া (māyā)}]{\bn{মায়া} (māyā) — tender attachment; fond pity}

\label{lesson:BN-C449-maya}



\begin{quote}
Before the new one: say the Bengali for to cry out in despair, then the Bengali for affection.
\end{quote}

\subsection*{You'll want to know: \bn{মায়া}}
\textbf{\bn{মায়া}} — \emph{māyā} — \textquotedblleft{}tender attachment; fond pity\textquotedblright{}.

\begin{quote}
\bn{পুরনো} \bn{বাড়িটা} \bn{ছাড়তে} \bn{আমার} \bn{মায়া} \bn{লাগল।}
\end{quote}

\begin{quote}
\emph{purôno bāṛiṭā chhāṛte āmār māyā lāglô.}
\end{quote}

\begin{quote}
Leaving the old house, I felt a pang of attachment.
\end{quote}

\bn{মায়া} is the soft pull you feel toward a person, a pet or a place; the older sense, illusion, belongs to philosophy.

\begin{cousinweb}
\textbf{\bn{মায়া}} (\emph{māyā}) is a tatsama, Sanskrit \emph{māyā}, \textquotedblleft{}supernatural power, illusion\textquotedblright{} (CDIAL 10060), where CDIAL, like Monier-Williams, also records a sense of compassion, from the lexicographers; the everyday Bengali sense of tender attachment probably grew from that one.
\end{cousinweb}

\begin{grammarlens}[title={What you've built}]
One more word for describing what you lived and felt.
\end{grammarlens}

\subsection*{Your turn}
\begin{itemize}
\raggedright
  \item \emph{Say it:} \emph{māyā}
  \item \emph{Say it:} \emph{māyā}, once more
  \item \emph{Say it:} \emph{purôno bāṛiṭā chhāṛte āmār māyā lāglô.} — Leaving the old house, I felt a pang of attachment.
  \item \emph{Recall:} say the Bengali for to cry out in despair, then the Bengali for affection, then say \emph{māyā} again
\end{itemize}

\begin{tcolorbox}[breakable,colback=teal!4,colframe=teal!35!black,title={Before you move on}]
What does \bn{মায়া} mean? (\textquotedblleft{}Tender attachment; fond pity\textquotedblright{}.) Say it once more.
\end{tcolorbox}
\section[\texorpdfstring{\bn{আদর} (ādôr)}{আদর (ādôr)}]{\bn{আদর} (ādôr) — affectionate care; cuddling and petting}

\label{lesson:BN-C449-ador}



\begin{quote}
Before the new one: say the Bengali for affection, then the Bengali for tender attachment.
\end{quote}

\subsection*{You'll want to know: \bn{আদর}}
\textbf{\bn{আদর}} — \emph{ādôr} — \textquotedblleft{}affectionate care; cuddling and petting\textquotedblright{}.

\begin{quote}
\bn{বাবা} \bn{মেয়েকে} \bn{আদর} \bn{করলেন।}
\end{quote}

\begin{quote}
\emph{bābā meyeke ādôr kôrlen.}
\end{quote}

\begin{quote}
Father gave his daughter a cuddle.
\end{quote}

\bn{আদর} \bn{করা}, to pet, to cuddle; \bn{আদরের}, beloved, the pet of the family.

\begin{cousinweb}
\textbf{\bn{আদর}} (\emph{ādôr}) is a tatsama, Sanskrit \emph{ādara}, \textquotedblleft{}respect, regard, care\textquotedblright{} (Monier-Williams); Bengali uses it for warm, petting affection, a sense the dictionary does not give.
\end{cousinweb}

\begin{grammarlens}[title={What you've built}]
One more word for describing what you lived and felt.
\end{grammarlens}

\subsection*{Your turn}
\begin{itemize}
\raggedright
  \item \emph{Say it:} \emph{ādôr}
  \item \emph{Say it:} \emph{ādôr}, once more
  \item \emph{Say it:} \emph{bābā meyeke ādôr kôrlen.} — Father gave his daughter a cuddle.
  \item \emph{Recall:} say the Bengali for affection, then the Bengali for tender attachment, then say \emph{ādôr} again
\end{itemize}

\begin{tcolorbox}[breakable,colback=teal!4,colframe=teal!35!black,title={Before you move on}]
What does \bn{আদর} mean? (\textquotedblleft{}Affectionate care; cuddling and petting\textquotedblright{}.) Say it once more.
\end{tcolorbox}
\section[\texorpdfstring{\bn{প্রেম} (prem)}{প্রেম (prem)}]{\bn{প্রেম} (prem) — love; romantic love}

\label{lesson:BN-C449-prem}



\begin{quote}
Before the new one: say the Bengali for tender attachment, then the Bengali for cuddling and petting.
\end{quote}

\subsection*{You'll want to know: \bn{প্রেম}}
\textbf{\bn{প্রেম}} — \emph{prem} — \textquotedblleft{}love; romantic love\textquotedblright{}.

\begin{quote}
\bn{কলেজে} \bn{পড়ার} \bn{সময়} \bn{অরুণ} \bn{আর} \bn{মীরার} \bn{প্রেম} \bn{হল।}
\end{quote}

\begin{quote}
\emph{kôleje pôṛār sômôy ôrun ār mīrār prem hôlô.}
\end{quote}

\begin{quote}
Arun and Mira fell in love while they were at college.
\end{quote}

\bn{প্রেমে} \bn{পড়া}, to fall in love.

\begin{cousinweb}
\textbf{\bn{প্রেম}} (\emph{prem}) is a tatsama, Sanskrit \emph{preman}, \textquotedblleft{}love, affection, kindness\textquotedblright{} (Monier-Williams).
\end{cousinweb}

\begin{grammarlens}[title={What you've built}]
One more word for describing what you lived and felt.
\end{grammarlens}

\subsection*{Your turn}
\begin{itemize}
\raggedright
  \item \emph{Say it:} \emph{prem}
  \item \emph{Say it:} \emph{prem}, once more
  \item \emph{Say it:} \emph{kôleje pôṛār sômôy ôrun ār mīrār prem hôlô.} — Arun and Mira fell in love while they were at college.
  \item \emph{Recall:} say the Bengali for tender attachment, then the Bengali for cuddling and petting, then say \emph{prem} again
\end{itemize}

\begin{tcolorbox}[breakable,colback=teal!4,colframe=teal!35!black,title={Before you move on}]
What does \bn{প্রেম} mean? (\textquotedblleft{}Love; romantic love\textquotedblright{}.) Say it once more.
\end{tcolorbox}
\section[\texorpdfstring{\bn{ভক্তি} \bn{করা} (bhôkti kôrā)}{ভক্তি করা (bhôkti kôrā)}]{\bn{ভক্তি} \bn{করা} (bhôkti kôrā) — to be devoted to; to revere (an elder)}

\label{lesson:BN-C449-bhokti-kora}



\begin{quote}
Before the new one: say the Bengali for cuddling and petting, then the Bengali for romantic love.
\end{quote}

\subsection*{You'll want to know: \bn{ভক্তি} \bn{করা}}
\textbf{\bn{ভক্তি} \bn{করা}} — \emph{bhôkti kôrā} — \textquotedblleft{}to be devoted to; to revere (an elder)\textquotedblright{}.

\begin{quote}
\bn{আমার} \bn{ভাই} \bn{বাবাকে} \bn{খুব} \bn{ভক্তি} \bn{করে।}
\end{quote}

\begin{quote}
\emph{āmār bhāi bābāke khub bhôkti kôre.}
\end{quote}

\begin{quote}
My brother is deeply devoted to our father.
\end{quote}

Past (he/she …): \bn{তিনি} \bn{তাঁর} \bn{বাবাকে} \bn{ভক্তি} \bn{করতেন}, \emph{tini tā̃r bābāke bhôkti kôrten}.

\bn{স্নেহ} goes down, to someone younger; \bn{ভক্তি} goes up, to a parent, an elder or a god.

\begin{cousinweb}
\textbf{\bn{ভক্তি} \bn{করা}} (\emph{bhôkti kôrā}) puts \bn{ভক্তি}, the tatsama Sanskrit \emph{bhakti}, devotion and fond attachment (Monier-Williams), before \bn{করা}, to do. \textbf{\bn{ভক্ত}}, a devotee, is the related Sanskrit \emph{bhakta}.
\end{cousinweb}

\begin{grammarlens}[title={What you've built}]
One more word for describing what you lived and felt.
\end{grammarlens}

\subsection*{Your turn}
\begin{itemize}
\raggedright
  \item \emph{Say it:} \emph{bhôkti kôrā}
  \item \emph{Say it:} \emph{bhôkti kôrā}, once more
  \item \emph{Say it:} \emph{āmār bhāi bābāke khub bhôkti kôre.} — My brother is deeply devoted to our father.
  \item \emph{Recall:} say the Bengali for cuddling and petting, then the Bengali for romantic love, then say \emph{bhôkti kôrā} again
\end{itemize}

\begin{tcolorbox}[breakable,colback=teal!4,colframe=teal!35!black,title={Before you move on}]
What does \bn{ভক্তি} \bn{করা} mean? (\textquotedblleft{}To be devoted to; to revere (an elder)\textquotedblright{}.) Say it once more.
\end{tcolorbox}
